\section{Related literature in detail}
\label{app:literature}

This appendix gives the full positioning summarized in \Cref{sec:literature}: the public-choice
analysis of sanctions, endogenous political power and conflict, hysteresis, the closest
antecedent, and the analytical methods.

This paper connects the public-choice analysis of sanctions with work on
endogenous political power, dynamic commitment, and hysteresis. Its
distinctive object is a slow target-side political state that sanctions
change and both governments value.

\subsection{Public choice, sanctions, and duration}

The public-choice starting point is \citet{Kaempfer1988}, who explain
sanctions through the political pressure and effectiveness of organized
interest groups.
\citet{HalcoussisKaempferLowenberg2021} emphasize that sanctions may
retain domestic political functions after their direct economic
leverage has diminished, while \citet{Galtung1967} supplies the classic
rally mechanism and \citet{Eaton1992} study sanctions effectiveness
without an endogenous political state. We make this political feedback
dynamic and strategic: the target chooses compliance, the sender
chooses sanctions, and both choices are derived from continuation
values generated by the same political state.

The four-case classification is what separates this paper from the two lines it
builds on. \citet{Eaton1992} already have both parties discounting, with
possibly different discount factors, under complete information, and derive
the outcome from relative toughness: a party's patience together with its
own cost of bearing sanctions; in their once-and-for-all case a more patient
sender can even lengthen the delay before compliance. The payoff-relevant
state there is the rival's standing choice and there is no political state, so
patience changes what is extracted and how long a confrontation lasts. In the public-choice
line, \citet{Kaempfer1988}
let the sanction level clear a static political market among interest
groups, and \citet{KaempferLowenbergMertens2004} study a dictator's choice of
power, loyalty and repression under sanctions; here both governments'
choices are derived from continuation values over a political state that the
sanction itself moves. The four cases pose a question of their own: holding
non-discount primitives fixed, which
government's forward-looking incentive produces non-compliance, and which changes its
price. Our answer is that the target's forward-looking incentive decides between
accommodation and non-compliance, while the sender's decides between extracting
compliance and sustaining entrenchment.

The model characterizes how long sanctions last once imposed. Its qualitative prediction is
that, after a temporary sanction shock carries the state across the political pivot, sanctions do
not fall within the equilibrium basin. This focus complements empirical
work on sanction duration and termination
\citep{BolksAlSowayel2000,AttiaGrauvogelVonSoest2020}, including the
sender--target termination dataset of \citet{AttiaGrauvogel2023}, in which
gradual lifting is common.

Two empirical findings delimit the political setting the model describes.
\citet{AhnLudema2020} find that a targeted regime shields the firms it
values, absorbing the sanction's cost through procurement terms so that the
burden falls elsewhere; we keep the government's political objective, the
constituencies it protects, and aggregate welfare as separate objects for
that reason, and the interest-group polity of \Cref{app:fourgroup}
does so explicitly. \citet{EscribaFolchWright2010} find that sanctions
destabilize personalist rulers, who lose the external revenue that funds
patronage, while single-party and military regimes absorb them through
taxation, reallocated spending and repression. The target's positive state
valuation $v_{\mathrm{lo}}>0$ of \Cref{sec:existence:placable} describes the second kind of
setting, in which an incumbent can absorb the sanction and the political
state $x$ measures the power of its anti-concession constituencies.
The sender's stake $w_{\mathrm{lo}}>0$ is a separate restriction, on the sender's own
constituencies, motivated by the interest-group literature above.

\subsection{Political persistence and strategic investment}
\label{sec:mechanism:comparison}

The political state is closest in interpretation to de facto power in
\citet{AcemogluRobinson2008}: influence that formal institutions do not
allocate, which there is a per-period investment by the group that holds it
\citep[p.~272]{AcemogluRobinson2008} and here a slow stock moved by the
sender's sanctions. It is also close to models in which current choices shape
future political power: extending the franchise changes the future median
voter and asset redistribution can consolidate democracy
\citep{AcemogluRobinson2001}, current decisions determine the future
distribution of political power \citep{AcemogluEgorovSonin2012}, institutional
states chosen by current power holders are shifted by shocks
\citep{AcemogluEgorovSonin2015}, and policy shapes future power
\citep{BaiLagunoff2011}. The difference is that two governments, with separate
discount factors, act on the same slow target-side state. Two of these papers name, as
directions for future work, ``capital-like state variables so that some
subcomponents of the state have autonomous dynamics''
\citep[Section~VI]{AcemogluEgorovSonin2012} and relations between countries
\citep[concluding section]{AcemogluEgorovSonin2015}. The present model takes up a
case of both: a political state with its own law of motion that another
country's instrument moves. The model shares with
the commitment-problem literature \citep{Powell2006} the feature that current
actions shift future political power in the absence of commitment. There the
parties bargain over transfers and the settlement's enforcement fails; here each
government sets its own instrument, compliance or the sanction, against a
political state that the sanction moves. The model's
dual-basin mechanism is related to \citet{HasslerEtAl2003}. The political
state here is itself shaped by forward-looking strategic choices.

The closest conflict-theoretic benchmark is the information-centered
framework of \citet{AcemogluWolitzkyMistrust}, in which mistrust,
misperception and misunderstanding drive conflict in a security dilemma, its
repeated version, and a deterrence model. Two comparisons place the present
paper. The first is what carries incentives across dates. In their security
dilemma and its repeated version it is a belief about the opponent's type; the
repeated version assumes infinitely impatient players to rule out dynamic
incentives from strategic experimentation, and names forward-looking leadership
as a direction for future research. In the dynamic deterrence result they report
from \citet[Theorem~2]{Wolitzky2023} it is the threat of switching to a worse
equilibrium after a publicly observed signal, which requires sufficiently
patient players. Here it
is an observed, payoff-relevant political state under complete information
and perfect monitoring, and the results compare constructed equilibria across
the four cases to isolate what each government's
forward-looking incentive does. The second is the stage structure, and here the benchmark
shares our logic: in their deterrence model with a perfectly observed action
the first mover's hawkish action is the better reply to any fixed response,
and yet the anticipated response induces restraint; likewise, in the class under
the lifting condition, compliance is costly at a fixed sanction and the sender's anticipated
response makes accommodation the target's unique optimum in the one-period problem. The difference is
that the sender ranks the two diagonal outcomes the
other way, so the two governments do not share a preference
for the restrained outcome (\Cref{sec:model:conflict-structure}), and that
both internalize an observed political state which sanctions move. The
constructed forward-looking equilibrium reverses the accommodation of the one-period problem for
that reason. In
\citet{AcemogluWolitzky2014} conflict spirals end from inside: Bayesian
players with limited memory come to treat a long spiral as uninformative about
the opponent's type and experiment with cooperation. Here information is
complete, and what resumed equilibrium
play does after a temporary shock depends on the state the shock leaves
behind.

A separate literature shows that valuing the future can sustain conflict
under complete information. In \citet{SkaperdasSyropoulos1996} and
\citet{GarfinkelSkaperdas2000}, parties who would cooperate or settle in a
one-shot interaction arm or fight once today's outcome determines tomorrow's
resources. \citet{Powell2004} gives a condition on continuation security
values under which no equilibrium of a stochastic game is efficient; in
\citet{Powell2012} sufficiently patient factions fight in every period to
forestall an adverse shift in power, and in \citet{KraininWiseman2016} enough
patience makes war occur in every subgame perfect equilibrium.
\citet{Powell2013} makes the distribution of power a continuous state chosen
by the government. We share the premise that the one-shot interaction ends in
accommodation and the dynamic one need not. Three things differ. Those models
give the players a common discount factor, or none, and where the players do
differ in patience (the long-lived litigant facing one-period challengers of
\citet[Appendix~B]{McBrideSkaperdas2009}), only one party's patience is
varied. The closest antecedents are \citet{Eaton1992}, who classify outcomes
by the target's cost of sanctions and its discount factor across
Markov-perfect steady states, and \citet{Eaton1999}, who report that a more
patient sender is more likely to succeed while the target's patience works
against it; in both the sender is committed between reviews, which is why
their sanctions are never suffered in equilibrium. Feedback Stackelberg games
also have existence results organized by patience: \citet{BressanJiang2020}
under a far-sighted leader or a narrow-sighted follower, and
\citet{Bucarey2023} re-solved at three corners of the discount factors. What is new here is
the object: our sender re-optimizes every period
against the political state its own sanction has built, which is why
sanctions are imposed and persist. That greater patience can deepen a political distortion is
published: investment in de facto power rises with the common discount factor
\citep[Proposition~2]{AcemogluRobinson2008}, populist bias rises with the
politicians' common discount factor \citep[Proposition~12]{AcemogluEgorovSonin2013},
patience makes a Pareto-inefficient state stable
\citep[Example~1]{AcemogluEgorovSonin2012}, and sufficiently patient factions
fight \citep{Powell2012}. A common discount factor can also act
non-monotonically, as on whether a society experiments with institutions
\citep[Proposition~7]{AcemogluEgorovSonin2015}. What is new here is that the two
governments' discount factors are separate and do different things: more target
patience moves the constructed equilibrium from relief to non-compliance, while a
forward-looking sender, against a myopic one, raises the price of relief on one
side and the intensity of conflict on the other. Sanctions theory has a standing
objection. Under complete information
and divisible demands the target concedes at the threat stage, so imposed
sanctions are instances where bargaining failed or information was
incomplete \citep{Drezner2003}; the same conclusion is an implication of the
binary model of \citet{Eaton1992}. The answer is what is proved. At
identical primitives, under stated conditions and with complete information
throughout, the realized
sanction in the constructed stationary equilibria is zero in the two cases with
a myopic target and positive in the two with a forward-looking one. The reason
is that the sender has a direct benefit
from sanctioning and the sanction builds entrenchment that the
forward-looking target values more than the relief it forgoes. The premise
that fails is that sanctions are a deadweight loss to both. The state here
is political power within the target polity, moved by the coercive instrument
itself and valued by both incumbents. Under the
conditions of \Cref{prop:relief-participation}, neither incumbent gains from an unconditional temporary suspension of sanctions followed by
the same full-non-compliance equilibrium. The condition has
the form of the participation conditions of the conflict literature
\citep{GarfinkelSkaperdas2000,McBrideSkaperdas2009}, and it presupposes
transferable utility when read as a side payment; it concerns uncompensated
relief. Compliance-contingent relief lies inside the
class through the lifting condition (compliance ends the sanction) and resembles
negotiated agreements of the kind the 2015 Iran arrangement provided.
The relief of the proposition is unconditional and set exogenously. This places it closer to
\citet{AcemogluTicchiVindigni2010}, where a civilian elite keeps the army too
small to end a civil war, because an army strong enough to end it could stage a
coup once the war is over and the elite cannot commit not to downsize it.

Relative to classic hysteresis and lock-in models
\citep{Dixit1989,Krugman1991,Arthur1989}, the slow state is moved by optimizing
agents in all four settings. In those models persistence comes from a sunk
cost, an externality or an increasing return. Here two strategic
governments move the state, one of them choosing the state-moving instrument while
internalizing the other's response, so that persistence is an equilibrium
property of a game. The dual-basin picture itself is known: Krugman's real-root case, in
which initial conditions alone select the long-run outcome; the equilibrium
of \citet{HasslerEtAl2003} with two basins inside one equilibrium, built by
splicing two rules; and the stable-state mapping of
\citet{AcemogluEgorovSonin2012}, which sends each initial state to the state
it settles in. What the paper adds is the state space and the
instrument: the basins here are basins of a political state moved by a
coercive instrument that one government chooses against the other's response.

History operates through that inherited state while the equilibrium rule
is held fixed. The dual-basin result therefore differs from
expectations-driven equilibrium selection: it varies the initial state and
holds fixed the equilibrium on which the governments coordinate. The distinction
also separates the paper's path dependence from models in which several
Markov equilibria coexist.

\subsection{Closest antecedent and analytical method}

The political-entrenchment motivation is also related to the earlier
working paper by \citet{GongZhong2025}. That paper studies a
target-only dynamic program with a reduced-form sanctions regime, reports
comparative statics in the target's patience among other parameters, and
supplements the model with sanction-duration evidence. With one decision
maker it cannot ask what the sender's patience does. The present paper makes
the sender a forward-looking player with its own discount factor and solves the
two-government feedback Stackelberg game. It obtains the separate roles of the
two discount factors and the threshold discount factor under the nine tightness
conditions, which make the transition affine and yield a threshold
independent of the sender's patience. It adds the high-basin and
shock results, the welfare
analysis, and the formal verification. The papers share a public-choice
motivation and the partial-adjustment law as a leading example; they have
distinct decision problems, empirical content, and results.

Methodologically, the paper combines a public-choice microfoundation,
a feedback Stackelberg Bellman system, and order methods for equilibrium
policies and basin dynamics
\citep{basar1999,krawczyk_petkov_multistage,StokeyLucas1989,
Curtat1996,Bertsekas2012,topkis1998,MilgromShannon1994,
QuahStrulovici2009,Krishnamurthy2023}. Interval dominance delivers a
monotone greatest response in cases not covered by classical single
crossing. The analysis combines that direction result with a separate
sanction-level condition, because monotonicity of a decision rule alone
does not imply an upward political path.

\paragraph{Monotone iteration and extremal steady states.}
The order-theoretic step behind \Cref{prop:order-selection} is classical.
Iteration of an increasing map from a lower solution converges to the minimal
fixed point above it, and from an upper solution to the maximal fixed point below
it \citep[Theorem 6.1 and Corollary 6.2]{Amann1976}; the statement we use is the
one-dimensional case, where continuity on a compact interval supplies the
compactness his theorem assumes. The paper connects this classical step to a
forward-looking political equilibrium, derives the relevant invariance and
monotonicity conditions from equilibrium restrictions, and separates convergence
of the political state from convergence of the equilibrium sanction. Terminology
for the economic sense of stability follows \citet{StokeyLucas1989}.

The closest primitive-level comparison is \citet{HuangPouzoRodriguezClare2025},
who construct a locator function whose sign determines the direction of motion,
read steady states and their stability off it, and obtain convergence of every
path to a common limit when the locator is single-crossing from above. Their
setting is single-agent dynamic optimization, in which an envelope step produces
the locator. Carrying that step to a Stackelberg game with endogenous compliance
and a general transition is open, and it is the natural route to a
primitive-level strengthening.

% Moved from the main text, 2026-10-16 (structure plan S1): the comparison with
% information-based accounts of conflict dynamics, formerly Section 5.3.
\subsection{Conflict under complete information}
\label{sec:discussion:misperception}

\central{In this model, conflict is sustained by a political state that the
conflict moves and both governments value, under complete information and perfect
monitoring.}\footnote{The absence of cycles, persistence after a suspension of sanctions and
the comparison with history-dependent strategies are proved in
\Cref{app:supporting-results,app:moved-ak,app:proofs:history}.} In the information-based account, misperception and misunderstanding,
with mistrust, play central roles in the onset, continuation and cessation of
conflict \citep{AcemogluWolitzkyMistrust}. In \citet{AcemogluWolitzky2014},
spirals that misperceptions start end from inside, as Bayesian players with
limited memory come to treat a long spiral as uninformative.
\Cref{tab:misperception} compares the two accounts interpretively.

\begin{table}[htbp]
\centering\footnotesize
\caption{Two accounts of conflict dynamics. The left column paraphrases the
information-based account; the right column summarizes results of this paper at
their stated scope.}
\label{tab:misperception}
\begin{tabular}{@{}p{0.18\linewidth}p{0.36\linewidth}p{0.4\linewidth}@{}}
\toprule
 & Information-based account & Political-state mechanism (this paper) \\
\midrule
What links dates & A belief about the opponent's type, or a threat to switch
equilibria & An observed political state that the conflict moves \\
What sustains conflict & Mistrust, misperception, misunderstanding & A sender that
gains from sanctioning a non-compliant target, and a forward-looking target that values
the entrenchment sanctions build at least as much as the relief it forgoes
(\Cref{thm:regime-separation}); under the strong entrenchment condition threats cannot
coerce it (\Cref{sec:results:security}) \\
What ends it & Belief revision: a long spiral comes to look uninformative, and
players experiment with cooperation \citep{AcemogluWolitzky2014} & A temporary
suspension of sanctions after which the resumed sanction rule stays at or below the state
(\Cref{sec:discussion:persistence}); in
the constructed equilibria of the four cases, no sanction arises when the target
does not value the future state (\Cref{thm:four-cells}) \\
Cycles & Recurring spirals, started by misperceptions, under limited memory
\citep{AcemogluWolitzky2014}; traps or eventual peace under other information
structures \citep[Section~4]{AcemogluWolitzkyMistrust} & None along deterministic
equilibrium paths under a sanction rule that does not fall with the state \\
What would have to differ & Mistrust, misperception, misunderstanding, for instance
through mediation \citep[Section~4.4]{AcemogluWolitzkyMistrust} & The state, when
the resumed sanction rule keeps sanctions at or below it after a suspension of sanctions;
in the special case of \Cref{prop:collapse-threshold}, the target's patience and
valuation of the state; the sender's stake changes the sanction level under
non-compliance (\Cref{sec:results:wedge}) \\
\bottomrule
\end{tabular}
\end{table}

\central{The political state that sustains sanctions can also end them.} A
temporary suspension of sanctions lowers the state, and in any economy a resumed sanction
rule that stays at or below the state after the shock delivers durable
de-escalation (\Cref{sec:discussion:persistence}). No belief is revised: the
conflict ends because the state has moved.

\central{Along deterministic equilibrium paths, the political-state mechanism also
produces no cycles under a sanction rule that does not fall with the state.} State
and sanction paths move in one direction and converge, and a path that returns to
a state stays there. The result covers the constructed equilibria of the four cases and
every equilibrium in $\mathcal E$ whose sanction rule is non-decreasing, on any
economy. Within the model, a cycle along an equilibrium path
needs a sanction rule that falls with the state somewhere. The cycles of
\citet{AcemogluWolitzky2014}, spirals that misperceived actions start, end and
restart, come from noise and limited memory.

\central{Read as comparative statics, the results say what would have to differ
for sanctions to end under the political-state mechanism.} With information
already complete, the levers are the state and the target's incentives. Under the
strong entrenchment condition, threats cannot coerce the target, since relief is the
harshest treatment the sender can impose (\Cref{sec:results:security}). A temporary
suspension of sanctions has a lasting effect if it leaves the state in a basin where resumed
equilibrium play keeps sanctions at zero. It has no
effect on the long-run limit where the equilibrium law of motion has an explicit
Lipschitz constant below one.
In the special case of \Cref{prop:collapse-threshold}, the target's patience, with
its valuation of the state, its cost and its sanction pain, decides whether a
full-non-compliance equilibrium exists. The sender's stake raises the sanction under
non-compliance above the myopic sanction (\Cref{sec:results:wedge}), and non-compliance persists
even with a myopic sender (case $10$ of \Cref{thm:four-cells}).

\subsubsection*{Timing and information in the sanctions game}
\label{sec:model:conflict-structure}

The game is a sequential game of complete information. Each period the sender
observes the target's compliance before responding, and both observe the political
state. Two questions relate it to models of conflict, the first being current
action incentives. On the class, at any fixed sanction more compliance strictly
lowers the target's payoff, so in payoffs non-compliance is its better reply, as the
hawkish action is for the first mover of a deterrence game
\citep{AcemogluWolitzkyMistrust}. What pays for compliance is the sender's
response. In the myopic equilibrium the target's unique optimum is the least
compliance that ends the myopic sanction (\Cref{prop:myopic-benchmark}). On the
quadratic--linear family (\Cref{sec:model:examples}), the target strictly prefers this
accommodation to
non-compliance met by the myopic sanction, while the sender ranks those two outcomes the
other way.
\central{On the family, the two governments therefore have opposed interests over
accommodation.}

\central{The second question is what carries incentives across dates: the
political state that the instrument itself moves.} In the constructed
forward-looking equilibrium the target's non-compliance and the sender's sanction both
maximize objectives that include the state's continuation value, so a credible
current incentive to accommodate coexists with dynamic non-compliance. The governments'
only actions are compliance and the sanction. Against each stationary equilibrium
studied here, no pure history-dependent deviation is profitable, nor is any finite
mixture of them.
